% cubic-ramp.tex

\section{Hypersonic flow over a concave surface.}
\label{sec:cubic-ramp}
%
This is one of the two hypersonic flows studied by Mohammadian\,\cite{mohammadian_1972a}
in the Imperial College gun tunnel, more than 4 decades ago.
The gun tunnel was operated with a total pressure of 1600\,psia (11.03\,MPa) and 
a total temperature of 1300\,K with a Mach 12.25 contoured nozzle.

\medskip
Since we don't have full information on the tunnel and its operating condition, 
we have to make a few assumptions.
First, we assume that the nozzle produced a parallel and uniform flow of ideal air 
with free stream conditions $p$ = 63.43\,Pa, $T$=41.92\,K and $u$ = 1.59\,km/s.
Using Sutherland's expression for the viscosity of air, 
we estimate the viscosity of the free stream to be $\mu = 2.593 \times 10^{-6}$\,Pa.s.
Taking a length scale $L = 1$\,inch, we compute a Reynolds number $Re_L = 0.86 \times 10^{5}$,
which is the same as that reported in the original paper\,\cite{mohammadian_1972a}.

\bigskip
\subsection{Input script (.py)}
%
Figure\,\ref{fig:cubic-ramp-geometry} shows the flow region, as modelled for simulation.
The region is very simple but we have divided it into 22 blocks so that the computational load 
can be shared across a number of CPU cores.

\begin{figure}[htbp]
 \centering
 \includegraphics[width=\textwidth,viewport=0 0 520 214,clip=true]{../2D/cubic-ramp/cubic-ramp-edited.pdf}
 \caption{Schematic view of the simulated flow region for the hypersonic flow
          over Mohammadian's concave ramp.}
 \label{fig:cubic-ramp-geometry}
\end{figure}

\medskip
The concave ramp is defined by
$$
y = \frac{1}{150} \, x^3
$$
where $x$ is in inches and the surface angle is less than 28$^o$.
Beyond that point, the ramp surface is assumed straight until the corner
at $x$ = 6.4\,inches.
The ramp surface temperature was assumed to be a constant 296\,K.

\noindent
\topbar
\lstinputlisting[language={}]{../2D/cubic-ramp/cubic-ramp.py}
\bottombar


\bigskip
\subsection{Running the simulation}
%
In terms of required computer time, this simulation is fairly demanding,
taking more than 12 hours on a 4-core workstation.
The job scipts submitted to the batch system are shown below.
Note that the preparation script sets up the mapping of the full set of 22 blocks to fit onto 4 MPI tasks.

\noindent
\topbar
\lstinputlisting[language={}]{../2D/cubic-ramp/prep.sh}\index{mpimap!example of use}
\bottombar\\
\topbar
\lstinputlisting[language={}]{../2D/cubic-ramp/run.sh}
\bottombar\\

\subsection{Results}
%
Figure\,\ref{fig:cubic-ramp-field-data} shows some of the flow field data 
at $t$=1\,ms after flow start.
This is sufficient time for the flow to reach steady state.

\begin{figure}[htbp]
 \centering
 \subfloat[Pressure field.]{\includegraphics[width=0.7\textwidth]
    {../2D/cubic-ramp/cubic-ramp-factor-2-log10p-field-1-ms.png}\label{fig:cubic-ramp-pressure}}\\
 \subfloat[Temperature field.]{\includegraphics[width=0.7\textwidth]
    {../2D/cubic-ramp/cubic-ramp-factor-2-T-field-1-ms.png}\label{fig:cubic-ramp-temperature}}\\
 \subfloat[Mach number.]{\includegraphics[width=0.7\textwidth]
    {../2D/cubic-ramp/cubic-ramp-factor-2-M-field-1-ms.png}\label{fig:cubic-ramp-mach}}
 \caption{Computed flow field at $t$=1\,ms.}
 \label{fig:cubic-ramp-field-data}
\end{figure}

\medskip
The logarithmic pressure field shows the leading-edge interaction shock propagating into the free-stream
flow but then intersecting the ramp about two-thirds of the way along its length.
The gradually increasing pressure can also be seen in the near-surface region, leading up to this intersection.

\medskip
The temperature field shows more clearly the boundary layer that grows along the concave surface.
It also shows the compression of the boundary-layer flow as the surface curves up.
The oblique shock propagating up over the top of the ramp is a combination of the reflected lead-edge
interaction shock and the coalesced compression waves from the curved ramp.

\clearpage

\medskip
Although the computed flow field looks plausible,
the real proof of success of the simulation is in comparison with the experimental data.
Figure\,\ref{fig:cubic-ramp-data-compare} shows the pressure and heat-transfer
along the surface of the ramp.
The simulation has done a good job of estimating the pressure distribution for the early part of the ramp,
and does a fair job all the way up to the sharp corner at the top of the ramp,
although there appears to be a mismatch in the location of this corner 
(indicated by the sharp drop in pressure at the right end of each data set).

\medskip
The simulation has also done a good job on the heat transfer estimate,
which has been computed from the field data using the script in Section\,\ref{cubic-ramp-post-processing}.
Mohammadian has not provided dimensional data in the original paper so we have normalized the simulation
data is the same way as can be best guessed (with the assistance of Andrew Knisely, University of Illinois).
First, we compute local heat transfer from the normal gradient of temperature at the ramp surface,
followed by Stanton number as
$$
St = \frac{q}{\rho U C_p (T_0 - T_w)}
$$
and then plot the combination $St . Re_x^{\frac{1}{2}}$, to remove the singularity in heat transfer
at the leading edge of the ramp.
Agreement for the early part of the plate to 110\,mm is excellent, so we can have some 
confidence in the codes ability to model hypersonic viscous interactions, however,
what happens beyond that point is not captured by this purely two-dimensional simulation.

\begin{figure}[htbp]
 \centering
 \subfloat[Pressure.]{\includegraphics[width=0.5\textwidth]
    {../2D/cubic-ramp/surface-pressure.pdf}\label{fig:cubic-ramp-surface-pressure}}
 \subfloat[Heat transfer.]{\includegraphics[width=0.5\textwidth]
    {../2D/cubic-ramp/surface-heat-transfer.pdf}\label{fig:cubic-ramp-heat-transfer}}
 \caption{Distribution of pressure and heat transfer along the concave ramp.
   Simulation data is recorded at $t$=1\,ms into the simulation.
   Experimental data is from Ref.\,\cite{mohammadian_1972a}.}
 \label{fig:cubic-ramp-data-compare}
\end{figure}

\medskip
The boundary-layer thickness, as identified by the outer edge of the maximum density gradient 
seen in a schlieren image was also provided by Mohammadian.
Figure\,\ref{fig:cubic-ramp-grad-rho-field} shows the density gradient field of the simulated
flow and, from this image, corresponding points were measured manually
(with the assistance of the \verb!g3data! program).
These boundary-layer thickness data are shown in Figure\,\ref{fig:cubic-ramp-delta-star}, 
together with the experimental values from Ref.\cite{mohammadian_1972a}.
Although there is some scatter in the simulation-derived data, comparison is good.

\begin{figure}[htb]
 \centering
 \includegraphics[width=0.7\textwidth]{../2D/cubic-ramp/cubic-ramp-factor-2-grad-rho-field-1-ms-zoom-with-scale.png}
 \caption{Computed density gradient field.}
 \label{fig:cubic-ramp-grad-rho-field}
\end{figure}

\begin{figure}[htb]
 \centering
 \includegraphics[width=0.7\textwidth]{../2D/cubic-ramp/delta-star.pdf}
 \caption{Estimates of the boundary-layer displacement thickness along the ramp.}
 \label{fig:cubic-ramp-delta-star}
\end{figure}


\bigskip
\subsection{Postprocessing to get heat transfer}
\label{cubic-ramp-post-processing}
%
The scripts below use the functions imported from \verb!e3_flow.py!
at a slightly higher level than in the cone20 example.
The first extracts the data for the cell nearest to the ramp surface
and uses that data to compute the expected shear stress and heat transfer at the surface.

\noindent
\topbar
\lstinputlisting[language={}]{../2D/cubic-ramp/surface_properties.py}
\bottombar

\subsection{Notes}
\begin{itemize}
 \item Plotting was done with the following GNUPlot scripts.
 \lstinputlisting[language={}]{../2D/cubic-ramp/surface-pressure.gnuplot}
 \lstinputlisting[language={}]{../2D/cubic-ramp/surface-heat-transfer.gnuplot}
\end{itemize}
